% Ported from an earlier report section (one-time conversion; edit here).
\subsection{Follow-up: real decisions}\label{jb:sec:followup}

A reviewer of this report raised two objections. First, the constructed cases may be too easy and too
synthetic: he would rather see fewer cases with higher fidelity, taken from real agent runs. Second, prompt
caching is the most likely pushback against routing: \emph{``for longer chats you have to be confident in
savings for the smaller model.''} We ran one follow-up study for each. This section reports the first; the caching study is in \cref{vt:sec:survey}. Nothing in
\cref{jb:sec:background,jb:sec:method,jb:sec:validate,jb:sec:dev,jb:sec:holdout,jb:sec:failures,jb:sec:latency} was
re-measured or changed; the follow-ups are separate runs with their own evidence and their own
preregistration.

\begin{keyidea}
Two new questions, two new kinds of evidence. \textbf{Real decisions:} do the judges still help when the
decisions come from real agent sessions instead of hand-written cases? \textbf{Switching cost:} once a
conversation is long, does moving it to a cheaper model still save money, given that each model keeps its
own prompt cache?
\end{keyidea}

\subsubsection{Why real traces, and how the cases were built}\label{jb:sec:traces}

\paragraph{What a trace-derived case is} During earlier benchmark sessions, the bundle's in-session judge
faced thousands of real \emph{read-shortcut} decisions: ``should the agent just read this file, or think?''
A \term{trace-derived case} is one such decision point, rebuilt exactly as the judge saw it. We kept only
steps that the in-session judge had \emph{sent to the host model} (it fell back), because only then do we
know what the host did next. Steps the judge automated were excluded (\TrFastRouted\ of them): their
``outcome'' is the judge's own pick, so labelling them would be circular.

\paragraph{Rebuilding and fingerprinting} The request bytes on the wire were never logged, so each request
was rebuilt from the session transcript. A rebuilt case was eligible only if it reproduced five logged
state numbers plus the logged hashes of the candidate ids and their order; \TrRebuildFail\ decisions failed
and were dropped. This is a fingerprint match, not a capture. From \TrMined\ eligible, de-duplicated
decisions we drew \TrDrawn\ (seed \TrSeed), stratified by suite and ``hard first'' using outcome-side
features only (more candidates, line-window candidates, clipped observations), never whether a logged judge
had been wrong. Whole tasks were assigned to holdout or dev before anyone labelled anything.

\paragraph{Labels} Two reviewers labelled every case blind: they saw the judge's view and the host's next
one or two actions, but not the outcome-derived label or any judge's answer. They agreed with each other
at Cohen's $\kappa = \TrKappaAB$, and with the outcome label at only $\kappa = \TrKappaOutcome$. The final
label is the reviewers' shared label: \TrAgreed\ cases matched the outcome label, \TrOverridden\ overrode
it, and the \TrDropped\ cases where the reviewers disagreed were dropped. The holdout has \TrN\ cases
(\TrReasonLabels\ labelled \code{reason}, \TrReadCases\ labelled with a read); the dev split has \TrDevN\
and is a screen.

\begin{keyidea}
\textbf{``What the host did next'' is not the same as ``what was enough.''} The host model might read a file
it did not need, or read a different, equally good file. A careful reviewer asks a stricter question: from
what the judge could see, was one prepared read clearly sufficient? The two disagree often (the low $\kappa$
above), so the reviewers' label is the primary one and the host's next read is kept as a secondary check.
\end{keyidea}

\paragraph{Replayed in each judge's own form} Jev received the bundle's own System One request; Laya and
the Qwen3 models went through the bundle's own local backends. GPT-6 Luna and Sol have no native form and
got the benchmark's choice question, as did nimble 9B and the tev1 models, because Ollama's System One
endpoint rejects the bundle's object-valued criteria (\cref{vt:sec:since}). The ``form'' column of
\cref{jb:tab:trace} records which.

\paragraph{Preregistration} The case hash, endpoints and decision rules were committed before the holdout
ran, and the runner refuses a holdout run unless they match. Two primary endpoints: the wrong-automatic rate
on all \TrN\ cases, and the number of \emph{correct automatic reads} among the \TrReadCases\ cases where a
read was enough (the savings a judge actually delivers). A new rule asks whether the read shortcut is worth
running at all: a judge is ``useful here'' if it makes at least \TrUsefulMinReads\ correct automatic reads
with a wrong-automatic upper 95\,\% bound below \TrUsefulMaxUpperPct\,\%. One erratum: a
\TrSmokeCases-case smoke test of the adapters on dev cases ran about \TrSmokeMinutes\ minutes before the
preregistration commit, although the file says it was committed before any judge saw a trace case. No
holdout case was involved. API spend for the whole study was \$\TrSpendTotal.

\subsubsection{Results on real decisions \texorpdfstring{\prereg}{[preregistered]}}\label{jb:sec:traceresults}

\begin{table}[tbp]
\centering\small
\caption{Real read-shortcut decisions, holdout \prereg: \TrN\ cases, \TrReps\ repetitions, the bundle's
read-shortcut gate. Per-case majorities over repetitions. ``Upper'' is the upper end of the 95\,\% Wilson
interval of the wrong-automatic rate. ``Correct reads'' counts correct automatic reads among the
\TrReadCases\ cases whose label is a read. Latency pools all valid requests. The last row is the reference
that never acts.}
\label{jb:tab:trace}
{\footnotesize\setlength{\tabcolsep}{2pt}\begin{tabular*}{\textwidth}{@{\extracolsep{\fill}}l l r r r r r r r r@{}}
\toprule
 & & \multicolumn{2}{c}{Accuracy} & \multicolumn{2}{c}{Wrong automatic} & & Correct & \multicolumn{2}{c}{Latency (ms)} \\
\cmidrule(lr){3-4}\cmidrule(lr){5-6}\cmidrule(lr){9-10}
Judge & form & correct & 95\,\% CI & count & upper & automatic & reads (of \TrReadCases) & p50 & p95 \\
\midrule
\textbf{Jev 1.13} & native & \TrJevAccK/\TrN & \TrJevAccCI & \TrJevWaK & \TrJevWaUpper\,\% & \TrJevCovK & \TrJevReads & \TrJevPfifty & \TrJevPninetyfive \\
nimble 9B & adapted & \TrNimbleAccK/\TrN & \TrNimbleAccCI & \TrNimbleWaK & \TrNimbleWaUpper\,\% & \TrNimbleCovK & \TrNimbleReads & \TrNimblePfifty & \TrNimblePninetyfive \\
GPT-6 Luna & adapted & \TrLunaAccK/\TrN & \TrLunaAccCI & \TrLunaWaK & \TrLunaWaUpper\,\% & \TrLunaCovK & \TrLunaReads & \TrLunaPfifty & \TrLunaPninetyfive \\
GPT-6.1 Sol & adapted & \TrSolAccK/\TrN & \TrSolAccCI & \TrSolWaK & \TrSolWaUpper\,\% & \TrSolCovK & \TrSolReads & \TrSolPfifty & \TrSolPninetyfive \\
Qwen3 8B & native & \TrQwenEightAccK/\TrN & \TrQwenEightAccCI & \TrQwenEightWaK & \TrQwenEightWaUpper\,\% & \TrQwenEightCovK & \TrQwenEightReads & \TrQwenEightPfifty & \TrQwenEightPninetyfive \\
Qwen3 0.6B & native & \TrQwenPtSixAccK/\TrN & \TrQwenPtSixAccCI & \TrQwenPtSixWaK & \TrQwenPtSixWaUpper\,\% & \TrQwenPtSixCovK & \TrQwenPtSixReads & \TrQwenPtSixPfifty & \TrQwenPtSixPninetyfive \\
tev1 4B & adapted & \TrTevFourAccK/\TrN & \TrTevFourAccCI & \TrTevFourWaK & \TrTevFourWaUpper\,\% & \TrTevFourCovK & \TrTevFourReads & \TrTevFourPfifty & \TrTevFourPninetyfive \\
Laya base & native & \TrLayaAccK/\TrN & \TrLayaAccCI & \TrLayaWaK & \TrLayaWaUpper\,\% & \TrLayaCovK & \TrLayaReads & \TrLayaPfifty & \TrLayaPninetyfive \\
Qwen3 4B & native & \TrQwenFourAccK/\TrN & \TrQwenFourAccCI & \TrQwenFourWaK & \TrQwenFourWaUpper\,\% & \TrQwenFourCovK & \TrQwenFourReads & \TrQwenFourPfifty & \TrQwenFourPninetyfive \\
tev1 0.8B & adapted & \TrTevPtEightAccK/\TrN & \TrTevPtEightAccCI & \TrTevPtEightWaK & \TrTevPtEightWaUpper\,\% & \TrTevPtEightCovK & \TrTevPtEightReads & \TrTevPtEightPfifty & \TrTevPtEightPninetyfive \\
\midrule
\emph{Always fall back} & --- & \TrAlwaysK/\TrN & \TrAlwaysCI & 0 & --- & 0 & 0 & --- & --- \\
\bottomrule
\end{tabular*}
}
\end{table}
\howtoreadtable{A useful judge sits high in the ``correct reads'' column and low in the ``wrong automatic''
column. Accuracy is shown, but it rewards doing nothing: \TrReasonLabels\ of the \TrN\ labels are
\code{reason}, so a judge that never acts scores \TrAlwaysK/\TrN.}

\paragraph{No judge met the usefulness bar} \TrNUseful\ of the \NumBase\ judges qualified. \TrNGeMinReads\
judges took at least \TrUsefulMinReads\ correct reads, and every one of them also made at least one wrong
automatic read. With \TrN\ cases, an upper bound below \TrUsefulMaxUpperPct\,\% is only possible with zero
wrong reads, so the rule effectively asked for a judge that saves several steps and never misfires.

\paragraph{Jev captured the most savings, and also misfired} Jev took \TrJevReads\ of the \TrReadCases\
sufficient reads automatically and made \TrJevWaK\ wrong automatic reads (\TrJevWaPct\,\% of cases, upper
bound \TrJevWaUpper\,\%). GPT-6 Luna took \TrLunaReads\ with \TrLunaWaK\ wrong; GPT-6.1 Sol \TrSolReads\
with \TrSolWaK\ wrong, but \TrSolTimeouts\ of its \TrSolAnswers\ answers missed the \TrTimeoutMs\,ms timeout
(its median per repetition was \TrSolPfiftyRangeS\,s) and it cost \TrSolJevCostX$\times$ as much as Jev.
Ignoring the timeout, the same Sol answers would have taken \TrSolNoTimeoutReads\ reads with
\TrSolNoTimeoutWa\ wrong; that is hypothetical, because a slower judge also makes the host wait.

\paragraph{No difference among Jev, Luna and Sol was detected} Jev versus Luna: Holm-adjusted $p = \TrJevLunaAccPHolm$ on
accuracy and \TrJevLunaWaPHolm\ on wrong automatic reads; Luna versus Sol: \TrLunaSolAccPHolm\ and
\TrLunaSolWaPHolm. Jev versus Sol was missing from the run's contrast list; a post-hoc exact McNemar test
gives $p = \TrJevSolAccP$ (accuracy) and \TrJevSolWaP\ (wrong automatic). On the \TrReadCases\ read cases
alone, Jev's extra reads over Luna and Sol give $p = \TrJevLunaReadsP$ and \TrJevSolReadsP. No candidate
replaced Jev under the default-judge rule (\TrNReplaces\ did).

\paragraph{The local judges fail in two ways} Laya and tev1 0.8B never cleared the \MinP\ gate, so they
never acted. Qwen3 4B never gave an answer at all (\cref{vt:sec:since}): all \TrQwenFourAbstained\ of
its answers became \code{reason} at probability 1. Nimble 9B, Qwen3 0.6B and Qwen3 8B acted freely and were
often wrong (\TrNimbleWaK, \TrQwenPtSixWaK\ and \TrQwenEightWaK\ wrong automatic reads).



\paragraph{Why accuracy misleads here} Always falling back scores \TrAlwaysK/\TrN, more than any judge
that acts. But \TrAlwaysSThreeK\ of those \TrAlwaysK\ come from the \TrSThree\ SWE-bench cases, where the
judge's view had been clipped before the issue text (\cref{jb:sec:caveats}); on the other \TrSOneTwo\ cases
it scores \TrAlwaysSOneTwoK. The useful question is not ``how often is the judge right'' but ``how many
sufficient reads does it capture for each unnecessary one.''

\paragraph{What a wrong automatic read costs} Every candidate here is a read-only file read. A wrong
automatic read replaces the host's next step with a read it did not need: one lost step and some extra
context, and no irreversible effect. That is a very different failure from clicking \emph{Buy now}
(\cref{jb:sec:worked}). Its cost was assumed, not measured. Real states are also long, about
\TrJevInputTokens\ input tokens, so Jev cost \$\TrJevCost\ per million decisions here, \TrJevCostXFirst$\times$
its cost on the constructed holdout.

\paragraph{Against the host's own next read} Scored against what the host actually read next (\TrHostReadCases\
cases where it read something), Jev has \TrJevHostReads\ correct automatic reads and \TrJevHostWa\ wrong;
Luna \TrLunaHostReads\ and \TrLunaHostWa; Sol \TrSolHostReads\ and \TrSolHostWa. None of Jev's \TrJevWaK\
wrong automatic reads matched the host's next read (\TrJevWaMatchHost). This analysis was chosen after the
results were seen and is not preregistered.

\subsubsection{Read this before generalizing}\label{jb:sec:caveats}

\begin{caution}[title={Read this before generalizing}]
\begin{itemize}[leftmargin=1.1em]
  \item \textbf{Scope and selection.} These are \TrN\ decisions that some in-session judge had already
    \emph{declined} and sent to the host. They are not a sample of all decisions: about half are labelled
    ``a read was enough'' by design, against about \TrPoolReadPct\,\% of the full pool. Jev made
    \TrJevOwnWa\ wrong automatic reads on the \TrJevOwnN\ cases from its own sessions (cases its
    in-session self had declined) and \TrJevOtherWa\ on the \TrJevOtherN\ cases from other judges'
    sessions.
  \item \textbf{Label overrides.} Reviewers overrode the host-derived label in \TrOverridden\ of
    \TrKept\ kept cases, leaning strict. Of the \TrHoldOverrides\ holdout overrides, \TrOverReadToReason\
    turned a read into \code{reason}, \TrOverReasonToRead\ turned \code{reason} into a read and
    \TrOverReadToRead\ swapped one read for another. Reviewers saw the host's next actions; they were blind
    only to the proposed label and to judge answers.
  \item \textbf{Task clipping.} Every case ran at a \TrStateBudget-character state budget (the bundle
    default is \TrStateDefault). In the \TrSThree\ SWE-bench holdout cases the task text the judge saw ended
    after about \TrTaskSeen\ characters, inside the operating rules: the judge never saw the issue.
  \item \textbf{Qwen3 4B is an artifact, not caution.} Its zero wrong automatic reads come from a defect
    in the candidate path, not from the model's judgement.
  \item \textbf{Clustering and source.} The \TrN\ cases come from \TrGroups\ task groups; Jev's
    \TrJevWaK\ wrong reads fall in \TrJevWrongGroups. All cases come from this bundle's own benchmark
    sessions, not human-driven traffic, and the reviewers are Anthropic models.
  \item \textbf{Erratum.} A \TrSmokeCases-case dev smoke test ran before the preregistration commit
    (\cref{jb:sec:traces}); no holdout case was involved.
\end{itemize}
\end{caution}

